\documentclass[11pt,a4paper]{article}
\usepackage[utf8]{inputenc}
\usepackage[T1]{fontenc}
\usepackage[spanish,es-minimal,es-noshorthands]{babel}
\usepackage{lmodern}
\usepackage[sfdefault,lining]{FiraSans}
\usepackage[a4paper,top=2.3cm,bottom=2.2cm,left=1.8cm,right=1.8cm,headheight=1.3cm,headsep=0.4cm,footskip=1cm]{geometry}
\usepackage{amsmath,amssymb}
\usepackage{icomma}
\usepackage{graphicx}
\usepackage[table]{xcolor}
\usepackage{tikz}
\usepackage{fontawesome5}
\usepackage{ccicons}
\usepackage{fancyhdr}
\usepackage{lastpage}
\usepackage{tasks}
\usepackage{needspace}
\usepackage[most]{tcolorbox}
\usepackage{array}
\usepackage[hidelinks]{hyperref}
\AtBeginDocument{\hypersetup{pdftitle={Soluciones de la ficha de ejercicios de la unidad 1: Números racionales. 3.º de ESO},pdfsubject={Matemáticas 3.º ESO},pdfkeywords={números enteros, fracciones, números racionales, decimales, fracción generatriz}}}

\definecolor{navy}{HTML}{1B2748}
\definecolor{accent}{HTML}{E8603A}
\definecolor{teal}{HTML}{1F9E90}
\definecolor{soft}{HTML}{EEF1F8}
\definecolor{muted}{HTML}{5D6685}

\setlength{\parindent}{0pt}
\setlength{\parskip}{3pt}
\settasks{label=\alph*), label-format=\bfseries\color{navy}, label-width=1.3em, item-indent=1.8em, column-sep=0.6em, after-item-skip=2pt, before-skip=3pt, after-skip=4pt}

% ---------- capçalera i peu ----------
\newcommand{\llicencia}{\href{https://creativecommons.org/licenses/by-nc-sa/4.0/deed.es}{\ccbyncsa}\ \ Licencia Creative Commons Atribución-NoComercial-CompartirIgual 4.0 Internacional (CC BY-NC-SA 4.0)}
\pagestyle{fancy}
\fancyhf{}
\fancyhead[L]{\raisebox{-0.35\height}{\includegraphics[height=0.95cm]{logo-gva.png}}}
\fancyhead[C]{\small\color{muted}Soluciones de la ficha · Unidad 1: Números racionales · 3.º de ESO}
\fancyhead[R]{\raisebox{-0.35\height}{\includegraphics[height=0.95cm]{logo-vermellar.png}}}
\fancyfoot[L]{\scriptsize\color{muted}\llicencia}
\fancyfoot[R]{\small\color{muted}\thepage}
\renewcommand{\headrulewidth}{0.4pt}
\renewcommand{\headrule}{\color{muted!40}\hrule height \headrulewidth}
\fancypagestyle{primera}{\fancyhf{}\renewcommand{\headrulewidth}{0pt}
  \fancyfoot[L]{\scriptsize\color{muted}\llicencia}
  \fancyfoot[R]{\small\color{muted}\thepage}}

% ---------- nivells i exercicis ----------
\newcommand{\estrelles}[1]{{\color{accent}\ifcase#1\or\faStar\faStar[regular]\faStar[regular]\or\faStar\faStar\faStar[regular]\or\faStar\faStar\faStar\fi}}
\newcounter{exer}
\newcommand{\ex}[1]{\par\needspace{4\baselineskip}\medskip\stepcounter{exer}\noindent{\color{navy}\bfseries\theexer.}\ \ {\small\estrelles{#1}}\ \ }

\newcommand{\seccio}[2]{\par\needspace{14\baselineskip}\bigskip
  \begin{tcolorbox}[colback=navy,colframe=navy,coltext=white,arc=2mm,boxrule=0pt,left=3mm,right=3mm,top=1.2mm,bottom=1.2mm]
  {\large\bfseries #1\quad #2}
  \end{tcolorbox}\vspace{-2pt}}

% ---------- recta real ----------
\newcommand{\rectabuida}{\begin{tikzpicture}[x=0.42cm,baseline=0]
  \draw[-latex,muted] (-6.7,0)--(6.8,0);
  \foreach \i in {-6,...,6}{\draw[muted] (\i,0.12)--(\i,-0.12) node[below=-1pt,font=\scriptsize,text=muted]{$\i$};}
  \end{tikzpicture}}
% \rectaint{x1}{x2}{estil esquerra}{estil dreta}: estil = ple, buit, inf
\newcommand{\extrem}[2]{\ifx#2p\fill[accent] (#1,0) circle (0.13);\else\ifx#2b\filldraw[fill=white,draw=accent,line width=1.2pt] (#1,0) circle (0.13);\fi\fi}
\newcommand{\rectaint}[4]{\begin{tikzpicture}[x=0.42cm,baseline=0]
  \draw[-latex,muted] (-6.7,0)--(6.8,0);
  \foreach \i in {-6,...,6}{\draw[muted] (\i,0.12)--(\i,-0.12) node[below=-1pt,font=\scriptsize,text=muted]{$\i$};}
  \draw[accent,line width=3pt] (#1,0)--(#2,0);
  \ifx#3i\fill[accent] (#1-0.25,0)--(#1+0.15,0.2)--(#1+0.15,-0.2)--cycle;\else\extrem{#1}{#3}\fi
  \ifx#4i\fill[accent] (#2+0.25,0)--(#2-0.15,0.2)--(#2-0.15,-0.2)--cycle;\else\extrem{#2}{#4}\fi
  \end{tikzpicture}}

\newcolumntype{C}[1]{>{\centering\arraybackslash}p{#1}}
\newcolumntype{M}[1]{>{\centering\arraybackslash}m{#1}}

% \solt: requadre per al resultat
\newcommand{\res}[1]{\boxed{#1}}
\newcommand{\si}{{\color{teal}\boldmath$\times$}}
\newcommand{\rectapunts}{\begin{tikzpicture}[x=0.42cm,baseline=0]
  \draw[-latex,muted] (-6.7,0)--(6.8,0);
  \foreach \i in {-6,...,6}{\draw[muted] (\i,0.12)--(\i,-0.12) node[below=-1pt,font=\scriptsize,text=muted]{$\i$};}
  \foreach \i in {-5,-2,3}{\fill[accent] (\i,0) circle (0.14);}
  \foreach \i in {5,2,-3}{\fill[teal] (\i,0) circle (0.14);}
  \end{tikzpicture}}

\def\fmneed{14}
% =====================================================================
% Capçalera FentMates de les fitxes: un color per a cada curs.
% Els colors són els de la pàgina inicial (LEVELS a assets/js/materials.js);
% el logo ve de _eines/logo/pdf/ (python3 _eines/logo/genera.py).
%
% Ús, al final del preàmbul de la unitat:
%   % =====================================================================
% Capçalera FentMates de les fitxes: un color per a cada curs.
% Els colors són els de la pàgina inicial (LEVELS a assets/js/materials.js);
% el logo ve de _eines/logo/pdf/ (python3 _eines/logo/genera.py).
%
% Ús, al final del preàmbul de la unitat:
%   % =====================================================================
% Capçalera FentMates de les fitxes: un color per a cada curs.
% Els colors són els de la pàgina inicial (LEVELS a assets/js/materials.js);
% el logo ve de _eines/logo/pdf/ (python3 _eines/logo/genera.py).
%
% Ús, al final del preàmbul de la unitat:
%   \input{../comu/fentmates.tex}
%   \fmcurs{1eso}
% i, al començament del document:
%   \capcalera{Matemàtiques}{1r ESO}{Fitxa d'exercicis · Unitat 2}{Divisibilitat}
% Opcional abans de \input: \def\fmneed{14} (espai mínim abans de cada secció).
% =====================================================================
\makeatletter
\def\fm@logodir{../../logo/pdf}
\providecommand{\fmneed}{10}
\def\fm@cursnom{}

% clau del curs, color i color del text damunt del color
\newcommand{\fm@curs}[3]{\@namedef{fm@color@#1}{#2}\@namedef{fm@text@#1}{#3}}
\fm@curs{1eso}{E6B84A}{navy}
\fm@curs{2eso}{3B82F6}{white}
\fm@curs{3eso}{FF7A50}{white}
\fm@curs{4eso-a}{A78BFA}{white}
\fm@curs{4eso-b}{6D28D9}{white}
\fm@curs{1batx-mi}{9B1C1C}{white}
\fm@curs{1batx-ccss}{EF6B6B}{white}
\fm@curs{2batx-mii}{1E6B47}{white}
\fm@curs{2batx-ccss}{2BB3A3}{white}

\newcommand{\fmcurs}[1]{%
  \@ifundefined{fm@color@#1}{\PackageError{fentmates}{Curs desconegut: #1}{}}{}%
  \edef\fm@tmp{\noexpand\definecolor{curs}{HTML}{\@nameuse{fm@color@#1}}%
               \noexpand\colorlet{curstext}{\@nameuse{fm@text@#1}}}\fm@tmp
  \colorlet{cursfosc}{curs!72!black}%
  \def\fm@logo{\fm@logodir/vertical-#1.pdf}%
  \def\fm@marca{\fm@logodir/marca-#1.pdf}}

% ---------- capçalera de la primera pàgina ----------
% #1 matèria (petit), #2 curs, #3 tipus de document i unitat, #4 títol
\newcommand{\capcalera}[4]{%
  \gdef\fm@cursnom{#2}%
  \par\noindent
  \begin{tcolorbox}[colback=curs,colframe=curs,boxrule=0pt,arc=2.2mm,width=3.5cm,height=2.15cm,
      valign=center,left=3mm,right=2mm,top=1mm,bottom=1mm,nobeforeafter,box align=center]
    \color{curstext}{\footnotesize #1}\par\vspace{1pt}{\fontsize{19}{21}\selectfont\bfseries #2\par}%
  \end{tcolorbox}\hspace{5mm}%
  \begin{minipage}[c]{\dimexpr\linewidth-3.5cm-5mm-2.9cm\relax}
    {\color{cursfosc}\small\bfseries\MakeUppercase{#3}}\par\vspace{2pt}
    {\color{navy}\fontsize{20}{23}\selectfont\bfseries #4\par}
  \end{minipage}\hfill
  \begin{minipage}[c]{2.6cm}\raggedleft\includegraphics[width=2.6cm]{\fm@logo}\end{minipage}\par
  \vspace{2pt}}

% ---------- barres de secció del color del curs ----------
\renewcommand{\seccio}[2]{\par\needspace{\fmneed\baselineskip}\bigskip
  \begin{tcolorbox}[colback=curs,colframe=curs,coltext=curstext,arc=2mm,boxrule=0pt,left=3mm,right=3mm,top=1.2mm,bottom=1.2mm]
  {\large\bfseries #1\quad #2}
  \end{tcolorbox}\vspace{-2pt}}

% ---------- estrelles de dificultat del color del curs (a Batxillerat no se'n posen) ----------
\@ifundefined{estrelles}{}{%
  \renewcommand{\estrelles}[1]{{\color{curs}\ifcase#1\or\faStar\faStar[regular]\faStar[regular]\or\faStar\faStar\faStar[regular]\or\faStar\faStar\faStar\fi}}}

% ---------- capçalera i peu de les pàgines ----------
% primera pàgina: logos institucionals al peu
\fancypagestyle{primera}{\fancyhf{}\renewcommand{\headrulewidth}{0pt}
  \fancyfoot[L]{\raisebox{-0.32\height}{\includegraphics[height=0.78cm]{logo-gva.png}}\hspace{5mm}%
    \raisebox{-0.32\height}{\includegraphics[height=0.78cm]{logo-vermellar.png}}}
  \fancyfoot[R]{\scriptsize\color{muted}\href{https://creativecommons.org/licenses/by-nc-sa/4.0/deed.ca}{\ccbyncsa}\ \ CC BY-NC-SA 4.0 · fentmates.org\qquad\small\thepage}}
% la resta: el símbol a l'esquerra i el curs a la dreta, amb el títol al centre com abans
\fancyhead[L]{\raisebox{-0.3\height}{\includegraphics[height=0.62cm]{\fm@marca}}}
\fancyhead[R]{\small\bfseries\color{cursfosc}\fm@cursnom}
\renewcommand{\headrule}{\color{curs!45}\hrule height \headrulewidth}
\makeatother

%   \fmcurs{1eso}
% i, al començament del document:
%   \capcalera{Matemàtiques}{1r ESO}{Fitxa d'exercicis · Unitat 2}{Divisibilitat}
% Opcional abans de \input: \def\fmneed{14} (espai mínim abans de cada secció).
% =====================================================================
\makeatletter
\def\fm@logodir{../../logo/pdf}
\providecommand{\fmneed}{10}
\def\fm@cursnom{}

% clau del curs, color i color del text damunt del color
\newcommand{\fm@curs}[3]{\@namedef{fm@color@#1}{#2}\@namedef{fm@text@#1}{#3}}
\fm@curs{1eso}{E6B84A}{navy}
\fm@curs{2eso}{3B82F6}{white}
\fm@curs{3eso}{FF7A50}{white}
\fm@curs{4eso-a}{A78BFA}{white}
\fm@curs{4eso-b}{6D28D9}{white}
\fm@curs{1batx-mi}{9B1C1C}{white}
\fm@curs{1batx-ccss}{EF6B6B}{white}
\fm@curs{2batx-mii}{1E6B47}{white}
\fm@curs{2batx-ccss}{2BB3A3}{white}

\newcommand{\fmcurs}[1]{%
  \@ifundefined{fm@color@#1}{\PackageError{fentmates}{Curs desconegut: #1}{}}{}%
  \edef\fm@tmp{\noexpand\definecolor{curs}{HTML}{\@nameuse{fm@color@#1}}%
               \noexpand\colorlet{curstext}{\@nameuse{fm@text@#1}}}\fm@tmp
  \colorlet{cursfosc}{curs!72!black}%
  \def\fm@logo{\fm@logodir/vertical-#1.pdf}%
  \def\fm@marca{\fm@logodir/marca-#1.pdf}}

% ---------- capçalera de la primera pàgina ----------
% #1 matèria (petit), #2 curs, #3 tipus de document i unitat, #4 títol
\newcommand{\capcalera}[4]{%
  \gdef\fm@cursnom{#2}%
  \par\noindent
  \begin{tcolorbox}[colback=curs,colframe=curs,boxrule=0pt,arc=2.2mm,width=3.5cm,height=2.15cm,
      valign=center,left=3mm,right=2mm,top=1mm,bottom=1mm,nobeforeafter,box align=center]
    \color{curstext}{\footnotesize #1}\par\vspace{1pt}{\fontsize{19}{21}\selectfont\bfseries #2\par}%
  \end{tcolorbox}\hspace{5mm}%
  \begin{minipage}[c]{\dimexpr\linewidth-3.5cm-5mm-2.9cm\relax}
    {\color{cursfosc}\small\bfseries\MakeUppercase{#3}}\par\vspace{2pt}
    {\color{navy}\fontsize{20}{23}\selectfont\bfseries #4\par}
  \end{minipage}\hfill
  \begin{minipage}[c]{2.6cm}\raggedleft\includegraphics[width=2.6cm]{\fm@logo}\end{minipage}\par
  \vspace{2pt}}

% ---------- barres de secció del color del curs ----------
\renewcommand{\seccio}[2]{\par\needspace{\fmneed\baselineskip}\bigskip
  \begin{tcolorbox}[colback=curs,colframe=curs,coltext=curstext,arc=2mm,boxrule=0pt,left=3mm,right=3mm,top=1.2mm,bottom=1.2mm]
  {\large\bfseries #1\quad #2}
  \end{tcolorbox}\vspace{-2pt}}

% ---------- estrelles de dificultat del color del curs (a Batxillerat no se'n posen) ----------
\@ifundefined{estrelles}{}{%
  \renewcommand{\estrelles}[1]{{\color{curs}\ifcase#1\or\faStar\faStar[regular]\faStar[regular]\or\faStar\faStar\faStar[regular]\or\faStar\faStar\faStar\fi}}}

% ---------- capçalera i peu de les pàgines ----------
% primera pàgina: logos institucionals al peu
\fancypagestyle{primera}{\fancyhf{}\renewcommand{\headrulewidth}{0pt}
  \fancyfoot[L]{\raisebox{-0.32\height}{\includegraphics[height=0.78cm]{logo-gva.png}}\hspace{5mm}%
    \raisebox{-0.32\height}{\includegraphics[height=0.78cm]{logo-vermellar.png}}}
  \fancyfoot[R]{\scriptsize\color{muted}\href{https://creativecommons.org/licenses/by-nc-sa/4.0/deed.ca}{\ccbyncsa}\ \ CC BY-NC-SA 4.0 · fentmates.org\qquad\small\thepage}}
% la resta: el símbol a l'esquerra i el curs a la dreta, amb el títol al centre com abans
\fancyhead[L]{\raisebox{-0.3\height}{\includegraphics[height=0.62cm]{\fm@marca}}}
\fancyhead[R]{\small\bfseries\color{cursfosc}\fm@cursnom}
\renewcommand{\headrule}{\color{curs!45}\hrule height \headrulewidth}
\makeatother

%   \fmcurs{1eso}
% i, al començament del document:
%   \capcalera{Matemàtiques}{1r ESO}{Fitxa d'exercicis · Unitat 2}{Divisibilitat}
% Opcional abans de \input: \def\fmneed{14} (espai mínim abans de cada secció).
% =====================================================================
\makeatletter
\def\fm@logodir{../../logo/pdf}
\providecommand{\fmneed}{10}
\def\fm@cursnom{}

% clau del curs, color i color del text damunt del color
\newcommand{\fm@curs}[3]{\@namedef{fm@color@#1}{#2}\@namedef{fm@text@#1}{#3}}
\fm@curs{1eso}{E6B84A}{navy}
\fm@curs{2eso}{3B82F6}{white}
\fm@curs{3eso}{FF7A50}{white}
\fm@curs{4eso-a}{A78BFA}{white}
\fm@curs{4eso-b}{6D28D9}{white}
\fm@curs{1batx-mi}{9B1C1C}{white}
\fm@curs{1batx-ccss}{EF6B6B}{white}
\fm@curs{2batx-mii}{1E6B47}{white}
\fm@curs{2batx-ccss}{2BB3A3}{white}

\newcommand{\fmcurs}[1]{%
  \@ifundefined{fm@color@#1}{\PackageError{fentmates}{Curs desconegut: #1}{}}{}%
  \edef\fm@tmp{\noexpand\definecolor{curs}{HTML}{\@nameuse{fm@color@#1}}%
               \noexpand\colorlet{curstext}{\@nameuse{fm@text@#1}}}\fm@tmp
  \colorlet{cursfosc}{curs!72!black}%
  \def\fm@logo{\fm@logodir/vertical-#1.pdf}%
  \def\fm@marca{\fm@logodir/marca-#1.pdf}}

% ---------- capçalera de la primera pàgina ----------
% #1 matèria (petit), #2 curs, #3 tipus de document i unitat, #4 títol
\newcommand{\capcalera}[4]{%
  \gdef\fm@cursnom{#2}%
  \par\noindent
  \begin{tcolorbox}[colback=curs,colframe=curs,boxrule=0pt,arc=2.2mm,width=3.5cm,height=2.15cm,
      valign=center,left=3mm,right=2mm,top=1mm,bottom=1mm,nobeforeafter,box align=center]
    \color{curstext}{\footnotesize #1}\par\vspace{1pt}{\fontsize{19}{21}\selectfont\bfseries #2\par}%
  \end{tcolorbox}\hspace{5mm}%
  \begin{minipage}[c]{\dimexpr\linewidth-3.5cm-5mm-2.9cm\relax}
    {\color{cursfosc}\small\bfseries\MakeUppercase{#3}}\par\vspace{2pt}
    {\color{navy}\fontsize{20}{23}\selectfont\bfseries #4\par}
  \end{minipage}\hfill
  \begin{minipage}[c]{2.6cm}\raggedleft\includegraphics[width=2.6cm]{\fm@logo}\end{minipage}\par
  \vspace{2pt}}

% ---------- barres de secció del color del curs ----------
\renewcommand{\seccio}[2]{\par\needspace{\fmneed\baselineskip}\bigskip
  \begin{tcolorbox}[colback=curs,colframe=curs,coltext=curstext,arc=2mm,boxrule=0pt,left=3mm,right=3mm,top=1.2mm,bottom=1.2mm]
  {\large\bfseries #1\quad #2}
  \end{tcolorbox}\vspace{-2pt}}

% ---------- estrelles de dificultat del color del curs (a Batxillerat no se'n posen) ----------
\@ifundefined{estrelles}{}{%
  \renewcommand{\estrelles}[1]{{\color{curs}\ifcase#1\or\faStar\faStar[regular]\faStar[regular]\or\faStar\faStar\faStar[regular]\or\faStar\faStar\faStar\fi}}}

% ---------- capçalera i peu de les pàgines ----------
% primera pàgina: logos institucionals al peu
\fancypagestyle{primera}{\fancyhf{}\renewcommand{\headrulewidth}{0pt}
  \fancyfoot[L]{\raisebox{-0.32\height}{\includegraphics[height=0.78cm]{logo-gva.png}}\hspace{5mm}%
    \raisebox{-0.32\height}{\includegraphics[height=0.78cm]{logo-vermellar.png}}}
  \fancyfoot[R]{\scriptsize\color{muted}\href{https://creativecommons.org/licenses/by-nc-sa/4.0/deed.ca}{\ccbyncsa}\ \ CC BY-NC-SA 4.0 · fentmates.org\qquad\small\thepage}}
% la resta: el símbol a l'esquerra i el curs a la dreta, amb el títol al centre com abans
\fancyhead[L]{\raisebox{-0.3\height}{\includegraphics[height=0.62cm]{\fm@marca}}}
\fancyhead[R]{\small\bfseries\color{cursfosc}\fm@cursnom}
\renewcommand{\headrule}{\color{curs!45}\hrule height \headrulewidth}
\makeatother

\fmcurs{3eso}

\begin{document}
\thispagestyle{primera}

% ---------- encapçalament ----------
\capcalera{Matemáticas}{3.º ESO}{Soluciones de la ficha · Unidad 1}{Números racionales}
\vspace{8pt}

\begin{tcolorbox}[colback=soft,colframe=soft,arc=2mm,left=3mm,right=3mm,top=1.5mm,bottom=1.5mm]
\small Intenta resolver cada ejercicio antes de mirar su solución. Si no te sale el mismo resultado, repasa los pasos o consulta los apuntes interactivos de la unidad. Cuando hay varias respuestas posibles, se muestra una de ejemplo.
\end{tcolorbox}

% =====================================================
\seccio{1}{Conjuntos numéricos: naturales, enteros y racionales}

\ex{1}\par\smallskip
{\renewcommand{\arraystretch}{2.1}\centering
\begin{tabular}{|M{2.6cm}|M{1.3cm}|M{1.3cm}|M{1.3cm}|}
\hline
\rowcolor{soft}\textbf{Número} & $\mathbb{N}$ & $\mathbb{Z}$ & $\mathbb{Q}$ \\ \hline
$15$ & \si & \si & \si \\ \hline
$-8$ & & \si & \si \\ \hline
$\dfrac34$ & & & \si \\ \hline
$0$ & \si & \si & \si \\ \hline
$-\dfrac{12}{3}=-4$ & & \si & \si \\ \hline
$0,6$ & & & \si \\ \hline
$1,\overline{3}$ & & & \si \\ \hline
$\dfrac{20}{4}=5$ & \si & \si & \si \\ \hline
$-2,5$ & & & \si \\ \hline
\end{tabular}\par}

\ex{2}
\begin{tasks}(4)
  \task $-4\in\mathbb{Z}$
  \task $7\in\mathbb{N}$
  \task $-\frac72\in\mathbb{Q}$
  \task $\frac18\in\mathbb{Q}$
  \task $5\in\mathbb{N}$
  \task $-6\in\mathbb{Z}$
  \task $\frac{25}{9}\in\mathbb{Q}$
  \task $4\in\mathbb{N}$
\end{tasks}

\ex{2}
\begin{tasks}(2)
  \task Verdadera: todo entero $a$ se puede escribir como $\frac{a}{1}$.
  \task Falsa: por ejemplo, $\frac12$ no es entero.
  \task Falsa: es un entero negativo.
  \task Falsa: $2-5=-3$ no es natural.
  \task Verdadera: por ejemplo, $\frac{5}{12}$ (porque $\frac{4}{12}<\frac{5}{12}<\frac{6}{12}$).
  \task Verdadera: $a:b=\frac{a}{b}$ es una fracción.
\end{tasks}

\ex{3}
\begin{tasks}(1)
  \task $\frac14=\frac{15}{60}$ y $\frac13=\frac{20}{60}$. Por ejemplo: $\frac{16}{60}=\frac{4}{15}$, \ $\frac{17}{60}$, \ $\frac{18}{60}=\frac{3}{10}$.
  \task Hay $6$: \ $-3,\ -2,\ -1,\ 0,\ 1,\ 2$.
  \task Por ejemplo, $-4$ es entero pero no natural, y $\frac23$ es racional pero no entero.
\end{tasks}

% =====================================================
\seccio{2}{Números enteros: valor absoluto, opuesto y orden}

\ex{1}
\begin{tasks}(4)
  \task $12$
  \task $9$
  \task $7$
  \task $-15$
  \task $4$
  \task $-6$
  \task $-3$
  \task $0$
\end{tasks}

\ex{1}
\begin{tasks}(1)
  \task $-12<-8<-3<-1<0<2<5$
  \task $7>4>1>-2>-7>-10$
  \task En naranja los números y en verde sus opuestos ($5$, $2$ y $-3$):\par\smallskip \rectapunts
\end{tasks}

\ex{2}
\begin{tasks}(3)
  \task $-5<-2$
  \task $|-5|>|-2|$
  \task $-10<0$
  \task $\text{op}(3)<\text{op}(-3)$\ \ ($-3<3$)
  \task $-(-4)>-4$
  \task $|-7|>6$
\end{tasks}

\ex{2}
\begin{tasks}(2)
  \task $x=8$\ \ o\ \ $x=-8$
  \task $-3,\ -2,\ -1,\ 0,\ 1,\ 2,\ 3$
  \task $-2,\ -1$
  \task $-3,\ -2,\ -1,\ 0,\ 1$
\end{tasks}

\ex{3}
\begin{tasks}(1)
  \task Miércoles ($-5$), lunes ($-3$), viernes ($-1$), jueves ($0$) y martes ($2$).
  \task $2-(-5)=2+5=\res{7\ ^\circ\text{C}}$
  \task Miércoles: $|-5|=5$.
\end{tasks}

% =====================================================
\seccio{3}{Suma y resta de números enteros}

\ex{1}
\begin{tasks}(6)
  \task $3$
  \task $-15$
  \task $-7$
  \task $11$
  \task $-10$
  \task $0$
\end{tasks}

\ex{1}
\begin{tasks}(2)
  \task $8+3-5=\res{6}$
  \task $-4-7+10=\res{-1}$
  \task $6-9+2=\res{-1}$
  \task $15-20+5-3=\res{-3}$
\end{tasks}

\ex{2}
\begin{tasks}(2)
  \task $(8+5)-(3+12+1)=13-16=\res{-3}$
  \task $(10+7+11)-(15+9)=28-24=\res{4}$
  \task $(13+18)-(20+6+5)=31-31=\res{0}$
  \task $(4+2+10)-(7+8+3)=16-18=\res{-2}$
\end{tasks}

\ex{2}
\begin{tasks}(2)
  \task $7-(-6)+(-3)=7+6-3=\res{10}$
  \task $-(-3)-8+4=3-8+4=\res{-1}$
  \task $12-[5-(-4)]=12-9=\res{3}$
  \task $-6-[7+4]-1=-6-11-1=\res{-18}$
  \task $[8-2]-[-4-(-4)]=6-0=\res{6}$
\end{tasks}

\ex{3}
\begin{tasks}(1)
  \task $-120+45-30+70=-35$. \ Se encuentra a $\res{35\text{ m}}$ de profundidad.
  \task $-212-(-287)=-212+287=\res{75\text{ años}}$
  \task $-4+11-9=\res{-2\ ^\circ\text{C}}$
\end{tasks}

% =====================================================
\seccio{4}{Multiplicación y división de números enteros}

\ex{1}
\begin{tasks}(4)
  \task $-30$
  \task $36$
  \task $-56$
  \task $1$
  \task $8$
  \task $-6$
  \task $-9$
  \task $0$
\end{tasks}

\ex{1}
\begin{tasks}(4)
  \task $30$
  \task $24$
  \task $-40$
  \task $0$
\end{tasks}

\ex{2}
\begin{tasks}(4)
  \task $-7$
  \task $-36$
  \task $-7$
  \task $-5$
\end{tasks}

\ex{2}
\begin{tasks}(2)
  \task $12:(-3)=\res{-4}$
  \task $12\cdot(-2)=\res{-24}$
  \task $(-25):(-5)=\res{5}$
  \task $36:(-4)=\res{-9}$
\end{tasks}

\ex{3}
\begin{tasks}(1)
  \task $6$ y $-6$: \ $6\cdot(-6)=-36$ y $6+(-6)=0$.
  \task Por ejemplo: $3$ y $-4$;\ \ $-3$ y $4$;\ \ $2$ y $-6$ (también $-2$ y $6$, $1$ y $-12$, $-1$ y $12$).
  \task $(-1)^{2025}=-1$ (exponente impar) y $(-1)^{2026}=1$ (par): \ $(-1)\cdot1=\res{-1}$.
\end{tasks}

% =====================================================
\seccio{5}{Operaciones combinadas con números enteros}

\ex{1}
\begin{tasks}(2)
  \task $5-6=\res{-1}$
  \task $-8-(-3)=-8+3=\res{-5}$
  \task $-12-(-10)=-12+10=\res{-2}$
  \task $-4+18=\res{14}$
\end{tasks}

\ex{2}
\begin{tasks}(2)
  \task $6-2\cdot(-4)=6+8=\res{14}$
  \task $(-5):(-5)+3=1+3=\res{4}$
  \task $-3\cdot(-2)^{2}=-3\cdot4=\res{-12}$
  \task $(-6)\cdot2+4=-12+4=\res{-8}$
  \task $-8+9-(-1)=-8+9+1=\res{2}$
  \task $-16+16=\res{0}$
\end{tasks}

\ex{2}
\begin{tasks}(2)
  \task $15-[4-2\cdot(-5)]=15-[4+10]=15-14=\res{1}$
  \task $-3\cdot[4-7]+(-5)=-3\cdot(-3)-5=9-5=\res{4}$
  \task $[-18+20]:(-2)+5=2:(-2)+5=-1+5=\res{4}$
  \task $2\cdot\{5-[3-(-5)]\}=2\cdot\{5-8\}=2\cdot(-3)=\res{-6}$
\end{tasks}

\ex{3}
\begin{tasks}(1)
  \task $9-4\cdot[6-(-8):(-4)]+(-2)=9-4\cdot[6-2]-2=9-16-2=\res{-9}$
  \task $\{-2\cdot[3-1\cdot5]+9\}:(-13)=\{-2\cdot(-2)+9\}:(-13)=13:(-13)=\res{-1}$
  \task $25:25-[12-16]\cdot(-1)=1-(-4)\cdot(-1)=1-4=\res{-3}$
  \task $3\cdot(4-6):2+1=3\cdot(-2):2+1=-3+1=\res{-2}$.\ \ Para que sea $4$: \ $(3\cdot4-6):2+1=6:2+1=\res{4}$ \ (o bien $3\cdot(4-6:2)+1=4$).
\end{tasks}

% =====================================================
\seccio{6}{Fracciones y números racionales}

\ex{1}
\begin{tasks}(6)
  \task $\dfrac23$
  \task $\dfrac34$
  \task $\dfrac45$
  \task $-\dfrac57$
  \task $\dfrac{6}{11}$
  \task $\dfrac73$
\end{tasks}

\ex{1}
\begin{tasks}(2)
  \task Sí: $3\cdot20=5\cdot12=60$.
  \task Sí: $4\cdot15=6\cdot10=60$.
  \task No: $7\cdot32=224$ y $8\cdot21=168$.
  \task Sí: $(-2)\cdot(-27)=9\cdot6=54$.
\end{tasks}

\ex{2}
\begin{tasks}(4)
  \task $x=10$
  \task $x=9$
  \task $x=3$
  \task $x=12$
\end{tasks}

\ex{2}
\begin{tasks}(1)
  \task $\text{m.c.m.}=12$: \ $\frac{9}{12},\ \frac{10}{12},\ \frac{8}{12},\ \frac{7}{12}$\quad$\Rightarrow$\quad $\dfrac{7}{12}<\dfrac23<\dfrac34<\dfrac56$
  \task $\text{m.c.m.}=20$: \ $-\frac{10}{20},\ \frac{12}{20},\ -\frac{15}{20},\ 0,\ \frac{14}{20}$\quad$\Rightarrow$\quad $-\dfrac34<-\dfrac12<0<\dfrac35<\dfrac{7}{10}$
\end{tasks}

\ex{3}
\begin{tasks}(1)
  \task $\frac25=\frac{8}{20}$ y $\frac35=\frac{12}{20}$: \ $\frac{9}{20}$, $\frac{10}{20}$ o $\frac{11}{20}$.
  \task $\dfrac{25}{60}=\res{\dfrac{5}{12}}$ de hora.
  \task $\frac{3k}{7k}$ con $3k+7k=60\Rightarrow k=6$: \ $\res{\dfrac{18}{42}}$.
\end{tasks}

% =====================================================
\seccio{7}{Operaciones con fracciones}

\ex{1}
\begin{tasks}(3)
  \task $\dfrac57$
  \task $-\dfrac39=\res{-\dfrac13}$
  \task $\dfrac36+\dfrac26=\res{\dfrac56}$
  \task $\dfrac{9}{12}-\dfrac{10}{12}=\res{-\dfrac{1}{12}}$
  \task $\dfrac{10}{5}-\dfrac35=\res{\dfrac75}$
  \task $\dfrac{3+4-6}{18}=\res{\dfrac{1}{18}}$
\end{tasks}

\ex{1}
\begin{tasks}(3)
  \task $\dfrac{18}{12}=\res{\dfrac32}$
  \task $-\dfrac{15}{60}=\res{-\dfrac14}$
  \task $-\dfrac{28}{8}=\res{-\dfrac72}$
  \task $\dfrac{30}{45}=\res{\dfrac23}$
  \task $\dfrac{84}{14}=\res{6}$
  \task $\dfrac{24}{3}=\res{8}$
\end{tasks}

\ex{2}
\begin{tasks}(2)
  \task $\dfrac{10+9-4}{24}=\dfrac{15}{24}=\res{\dfrac58}$
  \task $\dfrac{-4-9+5}{30}=-\dfrac{8}{30}=\res{-\dfrac{4}{15}}$
  \task $\dfrac{30-4-7}{10}=\res{\dfrac{19}{10}}$
  \task $\dfrac14-\left(-\dfrac16\right)=\dfrac{3}{12}+\dfrac{2}{12}=\res{\dfrac{5}{12}}$
\end{tasks}

\ex{2}
\begin{tasks}(2)
  \task $\dfrac{2\cdot7\cdot5^{2}}{3\cdot5\cdot2^{2}\cdot7}=\res{\dfrac56}$
  \task $-\dfrac{18\cdot21\cdot10}{35\cdot12\cdot9}=-\dfrac{3780}{3780}=\res{-1}$
  \task $\dfrac{12\cdot35}{25\cdot18}=\dfrac{2^{2}\cdot3\cdot5\cdot7}{5^{2}\cdot2\cdot3^{2}}=\res{\dfrac{14}{15}}$
  \task $\dfrac{8\cdot9}{27\cdot16}\cdot\dfrac32=\dfrac16\cdot\dfrac32=\res{\dfrac14}$
\end{tasks}

\ex{3}
\begin{tasks}(1)
  \task $\dfrac56-\dfrac38=\dfrac{20}{24}-\dfrac{9}{24}=\res{\dfrac{11}{24}}$
  \task $\dfrac23:\dfrac49=\dfrac{18}{12}=\res{\dfrac32}$
  \task $\dfrac{7}{12}-\left(-\dfrac14\right)=\dfrac{7}{12}+\dfrac{3}{12}=\dfrac{10}{12}=\res{\dfrac56}$
  \task El opuesto de $\frac27$ es $-\frac27$, y su inverso es $\res{-\dfrac72}$.
\end{tasks}

% =====================================================
\seccio{8}{Potencias de fracciones}

\ex{1}
\begin{tasks}(4)
  \task $\dfrac{9}{16}$
  \task $-\dfrac{8}{125}$
  \task $\dfrac{1}{81}$
  \task $-\dfrac{1}{16}$
  \task $1$
  \task $\dfrac94$
  \task $-4$
  \task $\dfrac{8}{27}$
\end{tasks}

\ex{2}
\begin{tasks}(3)
  \task $\left(\dfrac23\right)^{7}=\dfrac{128}{2187}$
  \task $\left(\dfrac54\right)^{4}=\dfrac{625}{256}$
  \task $\left(-\dfrac12\right)^{6}=\dfrac{1}{64}$
  \task $\left(\dfrac35\cdot\dfrac53\right)^{4}=1^{4}=1$
  \task $\left(\dfrac72\right)^{0}=1$
  \task $\left(\dfrac49:\dfrac23\right)^{3}=\left(\dfrac23\right)^{3}=\dfrac{8}{27}$
\end{tasks}

\ex{2}
\begin{tasks}(2)
  \task $\dfrac14+\dfrac18=\res{\dfrac38}$
  \task $\dfrac49-\dfrac19=\dfrac39=\res{\dfrac13}$
  \task $\dfrac94\cdot\dfrac{8}{27}=\dfrac{72}{108}=\res{\dfrac23}$
  \task $-\dfrac18:\dfrac{1}{16}=-\dfrac{16}{8}=\res{-2}$
\end{tasks}

\ex{3}
\begin{tasks}(2)
  \task $\left(\frac32\right)^{2}=\left(\frac23\right)^{-2}$:\ \ $\left(\dfrac23\right)^{5-2-2}=\res{\dfrac23}$
  \task $\left[\left(\dfrac12\right)^{-2}\right]^{2}=\left(\dfrac12\right)^{-4}=2^{4}=\res{16}$
  \task $\left(\dfrac35\right)^{6}:\left(\dfrac35\right)^{5}=\res{\dfrac35}$
  \task $\dfrac{16}{81}\cdot\left(-\dfrac{27}{8}\right)=-\dfrac{432}{648}=\res{-\dfrac23}$
\end{tasks}

% =====================================================
\seccio{9}{Operaciones combinadas con fracciones}

\ex{1}
\begin{tasks}(2)
  \task $\dfrac12+\dfrac14=\res{\dfrac34}$
  \task $\dfrac23-\dfrac{1}{12}=\dfrac{8}{12}-\dfrac{1}{12}=\res{\dfrac{7}{12}}$
  \task $\dfrac56\cdot\dfrac56=\res{\dfrac{25}{36}}$
  \task $\dfrac14:\dfrac58=\dfrac{8}{20}=\res{\dfrac25}$
\end{tasks}

\ex{2}
\begin{tasks}(2)
  \task $\dfrac23+\dfrac12\cdot\dfrac23=\dfrac23+\dfrac13=\res{1}$
  \task $\dfrac14\cdot\dfrac83=\dfrac{8}{12}=\res{\dfrac23}$
  \task $\dfrac35-\dfrac25:\dfrac16=\dfrac35-\dfrac{12}{5}=\res{-\dfrac95}$
  \task $\dfrac14+\dfrac34:\dfrac12=\dfrac14+\dfrac32=\res{\dfrac74}$
\end{tasks}

\ex{2}
\begin{tasks}(2)
  \task $\left[\dfrac23-\dfrac14\right]:\dfrac56=\dfrac{5}{12}:\dfrac56=\dfrac{30}{60}=\res{\dfrac12}$
  \task $1-\left[\dfrac34-\dfrac16\right]=1-\dfrac{7}{12}=\res{\dfrac{5}{12}}$
  \task $\dfrac49:\dfrac43-\dfrac16=\dfrac13-\dfrac16=\res{\dfrac16}$
  \task $\dfrac32-\dfrac23\cdot\left[\dfrac14-\dfrac34\right]=\dfrac32-\dfrac23\cdot\left(-\dfrac12\right)=\dfrac32+\dfrac13=\res{\dfrac{11}{6}}$
\end{tasks}

\ex{3}
\begin{tasks}(1)
  \task $\left\{\dfrac12-\left[\dfrac23-\dfrac{7}{12}\right]\right\}:\dfrac19=\left\{\dfrac12-\dfrac{1}{12}\right\}:\dfrac19=\dfrac{5}{12}\cdot9=\dfrac{45}{12}=\res{\dfrac{15}{4}}$
  \task $\left(\dfrac12\right)^{-2}\cdot\left(-\dfrac14\right)+\left(-\dfrac18\right)\cdot8=4\cdot\left(-\dfrac14\right)-1=-1-1=\res{-2}$
  \task $\dfrac{\frac16}{\frac43}=\dfrac{3}{24}=\res{\dfrac18}$
  \task $1+\frac12=\frac32$;\ \ $\frac{1}{\frac32}=\frac23$;\ \ $2-\frac23=\frac43$;\ \ $\frac12+\frac13=\frac56$:\ \ $\dfrac43:\dfrac56=\dfrac{24}{15}=\res{\dfrac85}$
\end{tasks}

% =====================================================
\seccio{10}{Problemas con números racionales}

\ex{1} $\dfrac{5}{12}\cdot840=\dfrac{5\cdot840}{12}=350$. \ Viven allí $\res{350\text{ alumnos}}$ y vienen de fuera $840-350=\res{490}$.

\ex{1} $\dfrac35\cdot165=99$. \ Ha hecho $\res{99\text{ km}}$ y le faltan $\res{66\text{ km}}$.

\ex{1} Una parte: $45:3=15$; \ la cantidad tiene $8$ partes: \ $15\cdot8=\res{120\text{ €}}$.

\ex{2} Novelas y ciencia: $\frac25+\frac14=\frac{13}{20}$. El resto es $\frac{7}{20}$ del total, que son $840$ libros:\\
$840:\dfrac{7}{20}=\dfrac{840\cdot20}{7}=\res{2\,400\text{ libros}}$. \ Novelas: $\dfrac25\cdot2\,400=\res{960}$.

\ex{2} Primero: $\frac13\cdot600=\res{200\text{ €}}$. \ Segundo: $\frac25\cdot600=\res{240\text{ €}}$. \ Tercero: $600-200-240=\res{160\text{ €}}$ \ (son los $\frac{4}{15}$).

\ex{2} Litros totales: $6\cdot\frac34=\frac92$ L. \ Vasos: $\dfrac92:\dfrac{3}{16}=\dfrac{144}{6}=\res{24\text{ vasos}}$.

\ex{2} Después del lunes queda $\frac23$. El martes se gasta $\frac34\cdot\frac23=\frac12$ del depósito.\\
Queda $\dfrac23-\dfrac12=\res{\dfrac16}$ del depósito, es decir, $\dfrac16\cdot1\,200=\res{200\text{ L}}$.

\ex{3} Después del alquiler queda $\frac35$ del sueldo. Comida: $\frac13\cdot\frac35=\frac15$ del sueldo.\\
Gasta $\frac25+\frac15=\frac35$; le queda $\frac25$ del sueldo, que son $600$ €: \ $600:\dfrac25=\res{1\,500\text{ €}}$.\\
Alquiler: $\frac25\cdot1\,500=600$ €. \ Comida: $\frac15\cdot1\,500=300$ €. \ (Comprobación: $1\,500-600-300=600$ \faCheck)

\ex{3}
\begin{tasks}(1)
  \task En una hora llenan $\frac14+\frac{1}{12}=\frac{3}{12}+\frac{1}{12}=\frac{4}{12}=\frac13$ del depósito $\Rightarrow$ tardan $\res{3\text{ horas}}$.
  \task En una hora: $\frac13-\frac16=\frac16$ del depósito $\Rightarrow$ tardan $\res{6\text{ horas}}$.
\end{tasks}

\ex{3} Primer día: $\frac38$; quedan $\frac58$. Segundo día: $\frac25\cdot\frac58=\frac14$. Vendidas: $\frac38+\frac14=\frac58$.\\
Quedan $\frac38$ de las sandías, que son $60$: \ $60:\dfrac38=\dfrac{60\cdot8}{3}=\res{160\text{ sandías}}$.

\ex{3} $81\cdot\left(\dfrac23\right)^{3}=81\cdot\dfrac{8}{27}=\res{24\text{ cm}}$ \ (los botes llegan a $54$, $36$ y $24$ cm).

% =====================================================
\seccio{11}{Números decimales: clasificación}

\ex{1}
\begin{tasks}(2)
  \task Exacto.
  \task Periódico puro $0,\overline{4}$: parte entera $0$, periodo $4$.
  \task Periódico mixto $1,2\overline{3}$: parte entera $1$, anteperiodo $2$, periodo $3$.
  \task No periódico (es $\pi$, un número irracional).
  \task Periódico puro $0,\overline{12}$: parte entera $0$, periodo $12$.
  \task Periódico mixto $5,1\overline{06}$: parte entera $5$, anteperiodo $1$, periodo $06$.
\end{tasks}

\ex{1}
\begin{tasks}(3)
  \task $0,875$ \ exacto
  \task $0,\overline{2}$ \ puro
  \task $0,8\overline{3}$ \ mixto
  \task $1,\overline{18}$ \ puro
  \task $0,35$ \ exacto
  \task $0,7\overline{3}$ \ mixto
\end{tasks}

\ex{2}
\begin{tasks}(2)
  \task $40=2^{3}\cdot5$ $\Rightarrow$ exacto.
  \task $21=3\cdot7$ $\Rightarrow$ periódico puro.
  \task $30=2\cdot3\cdot5$ $\Rightarrow$ periódico mixto.
  \task $\frac{9}{12}=\frac34$; \ $4=2^{2}$ $\Rightarrow$ exacto.
  \task $33=3\cdot11$ $\Rightarrow$ periódico puro.
  \task $18=2\cdot3^{2}$ $\Rightarrow$ periódico mixto.
  \task $\frac{14}{35}=\frac25$ $\Rightarrow$ exacto.
  \task $55=5\cdot11$ $\Rightarrow$ periódico mixto.
\end{tasks}

\ex{3}
\begin{tasks}(1)
  \task El numerador tiene que ser múltiplo de $3$ para que, al simplificar, desaparezca el $3$ del denominador ($12=2^{2}\cdot3$). Por ejemplo, $\frac{3}{12}=\frac14=0,25$ \ o \ $\frac{9}{12}=\frac34=0,75$.
  \task $0,3<0,\overline{3}<0,\overline{34}<0,3\overline{4}$ \ \ ($0,3000\ldots<0,3333\ldots<0,3434\ldots<0,3444\ldots$)
  \task Sí: $\frac13+\frac16=\frac26+\frac16=\frac36=\frac12=0,5$, que es exacto.
\end{tasks}

% =====================================================
\seccio{12}{Paso de decimal a fracción}

\ex{1}
\begin{tasks}(4)
  \task $\dfrac{8}{10}=\res{\dfrac45}$
  \task $\dfrac{245}{100}=\res{\dfrac{49}{20}}$
  \task $\dfrac{36}{1000}=\res{\dfrac{9}{250}}$
  \task $\dfrac{3125}{1000}=\res{\dfrac{25}{8}}$
\end{tasks}

\ex{2}
\begin{tasks}(2)
  \task $\res{\dfrac79}$
  \task $\dfrac{23-2}{9}=\dfrac{21}{9}=\res{\dfrac73}$
  \task $\dfrac{18}{99}=\res{\dfrac{2}{11}}$
  \task $\dfrac{136-1}{99}=\dfrac{135}{99}=\res{\dfrac{15}{11}}$
\end{tasks}

\ex{2}
\begin{tasks}(2)
  \task $\dfrac{16-1}{90}=\dfrac{15}{90}=\res{\dfrac16}$
  \task $\dfrac{127-12}{90}=\dfrac{115}{90}=\res{\dfrac{23}{18}}$
  \task $\dfrac{3-0}{90}=\res{\dfrac{1}{30}}$
  \task $\dfrac{2136-21}{990}=\dfrac{2115}{990}=\res{\dfrac{47}{22}}$
\end{tasks}

\ex{3}
\begin{tasks}(1)
  \task $N=0,999\ldots$: \ $10N=9,999\ldots$; restando, $9N=9$, es decir, $N=\res{1}$.
  \task $\dfrac39+\dfrac69=\dfrac13+\dfrac23=\res{1}$
  \task $0,\overline{45}=\frac{45}{99}=\frac{5}{11}=\frac{450}{990}$ \ y \ $0,4\overline{5}=\frac{45-4}{90}=\frac{41}{90}=\frac{451}{990}$. \ Es mayor $\res{0,4\overline{5}}$.
  \task $\dfrac{142857}{999999}=\res{\dfrac17}$
\end{tasks}

% =====================================================
\seccio{13}{Operaciones combinadas con fracciones y decimales}

\ex{2}
\begin{tasks}(2)
  \task $\dfrac12+\dfrac13=\res{\dfrac56}$
  \task $\dfrac65\cdot\dfrac23=\dfrac{12}{15}=\res{\dfrac45}$
  \task $\dfrac23-\dfrac16=\dfrac36=\res{\dfrac12}$
  \task $\dfrac73:\dfrac79=\dfrac{63}{21}=\res{3}$
\end{tasks}

\ex{2}
\begin{tasks}(2)
  \task $\dfrac14+\dfrac43\cdot\dfrac34=\dfrac14+1=\res{\dfrac54}$
  \task $\left(\dfrac13+\dfrac12\right)^{2}=\left(\dfrac56\right)^{2}=\res{\dfrac{25}{36}}$
  \task $\dfrac32-\dfrac29:\dfrac49=\dfrac32-\dfrac12=\res{1}$
  \task $\left(\dfrac{12}{5}-\dfrac23\right)\cdot\dfrac16=\dfrac{26}{15}\cdot\dfrac16=\dfrac{26}{90}=\res{\dfrac{13}{45}}$
\end{tasks}

\ex{3}
\begin{tasks}(1)
  \task $1,\overline{6}=\frac53$, \ $0,5=\frac12$, \ $0,\overline{3}=\frac13$: \ \ $\left(\dfrac53-\dfrac12\right):\left(\dfrac13\right)^{2}=\dfrac76:\dfrac19=\dfrac{63}{6}=\res{\dfrac{21}{2}}$
  \task $1,1\overline{6}=\frac{116-11}{90}=\frac{105}{90}=\frac76$: \ \ $\dfrac{\frac13+\frac14}{\frac76}=\dfrac{7}{12}:\dfrac76=\res{\dfrac12}$
  \task $0,\overline{9}=1$ \ y \ $1,\overline{1}=\frac{10}{9}$: \ \ $\left(1-\dfrac{1}{10}\right)\cdot\dfrac{10}{9}=\dfrac{9}{10}\cdot\dfrac{10}{9}=\res{1}$
  \task $2,0\overline{3}=\frac{203-20}{90}=\frac{183}{90}=\frac{61}{30}$: \ \ $\dfrac{61}{30}\cdot\dfrac35-\dfrac15=\dfrac{61}{50}-\dfrac{10}{50}=\res{\dfrac{51}{50}}$
\end{tasks}

\end{document}
